%% abstract.tex: abstract in PT and EN  (FEUP regulations)
%% -------------------------------------------------------

\chapter*{Abstract}
%
Demand forecasting is a foundational task in retail optimization that directly
influences supply chain coordination, inventory management, and promotional
strategies. While machine learning models have significantly advanced sequential
prediction, traditional approaches treat individual products as isolated time
series or attempt to learn shared behaviors implicitly, discarding explicit
relational dependencies such as substitution, complementarity, or shared
promotional impacts. Although Graph Neural Networks (GNNs) have shown notable
success in physical spatiotemporal domains, such as traffic and energy networks,
their application in retail remains severely limited owing to the latent, volatile,
and the time-varying nature of product-to-product relationships. To bridge this gap,
this dissertation presents a comprehensive spatiotemporal framework spanning two
integration paradigms: a decoupled, loss-independent embedding pipeline and an
integrated, end-to-end spatiotemporal architecture. In both, the product networks
are constructed identically through target-centric topologies over rolling
historical windows using data-driven, \emph{loss-independent} statistical
similarity and distance criteria, namely Spearman's correlation and
Complexity-Invariant Distance (CID). The two paradigms differ in how relational
features are then extracted: the decoupled strategy uses Graph2Vec to learn graph
embeddings without any access to the forecasting objective, whereas the end-to-end
strategy uses inductive Graph Convolutional Networks (GCN) and Graph Attention
Networks (GAT) whose parameters are optimized \emph{jointly} with the forecasting
loss. The resulting context vectors are fused with temporal inputs and evaluated
through recursive multi-step forecasting heads, including Long Short-Term Memory
(LSTM) networks and Time-Distributed Multilayer Perceptrons (MLP). The empirical
pipeline was evaluated on an operational retail dataset partitioned into
chronological training, validation, and holdout segments. Product prioritization is
guided by the ADI-$CV^2$ demand-pattern framework, which curates a high-signal
benchmark subset of 60 top-selling smooth and erratic items from a 1,427-product
catalog. The framework was validated using non-parametric statistical criteria
alongside point-forecasting metrics, namely MAE and POCID. The results reveal a
clear magnitude--direction trade-off: relational augmentation does not improve, and
can slightly degrade, absolute error (MAE), where univariate baselines and the
loss-independent Graph2Vec embeddings remain strongest, but it significantly and
robustly improves directional accuracy (POCID), with end-to-end attention-based GAT
hybrids best at anticipating synchronized demand movements. The work thus
establishes loss-independent graph-\emph{construction} guidelines and characterizes
when relational intelligence is, and is not, beneficial for local retail
forecasting.

Keywords: Retail Demand Forecasting, Graph Neural Networks, Spatiotemporal
Modeling, Dynamic Graph Embedding, Inductive Representation Learning.


\acresetall %% to reset the acronym usage
